% Topic document -- target: 2 to 4 pages, 11pt, 1-inch margins.
\documentclass[11pt]{article}
\usepackage[utf8]{inputenc}
\usepackage[T1]{fontenc}
\usepackage[margin=0.88in]{geometry}
\usepackage{amsmath,amssymb}
\usepackage{booktabs,array}
\usepackage{xcolor}
\usepackage[round]{natbib}
\usepackage[colorlinks=true,allcolors=blue!50!black]{hyperref}
\setlength{\bibsep}{1.5pt}
\setlength{\parskip}{2pt}

\title{Relative Loan Pricing as an Early-Warning Signal of Non-Performing Loans\\
\large Endogenous Screening under Borrower Opacity}
\author{Fabián Matías Galarza Chumbe\\
\small\url{https://github.com/fgalarzac/ai-project}}
\date{Topic presentation, October 14, 2026}

\begin{document}
\maketitle
\vspace{-1.2em}

\section{Question, motivation, and track}

A bank supervisor observes loan prices before portfolios mature. Can a bank's
unusually high price serve as an early-warning signal for its future
non-performing loan ratio (NPL ratio)? The sign is not obvious. A high price
can worsen the applicant pool through adverse selection, but it also makes a
performing loan more valuable and may induce the bank to screen more carefully.

\textbf{Question.} When does a higher loan spread relative to comparable banks
predict a higher future NPL ratio, and when can it predict a lower one? This is
a \textbf{Track B} project. My thesis proposed the two mechanisms in reduced
form; this adaptation adds the bank's missing optimization problem, derives
the screening response from primitives, and states the restrictions required
for the predictive sign to reverse with borrower opacity. The claim is about
forecasting, not a causal effect of forcing a bank to change its rate.

Index banks by \(j\), comparable credit segments by \(s\), pricing months by
\(t\), and forecast horizons by \(h\). The empirical signal is built in two
steps:
\begin{equation}
s_{jst}=i_{jst}-r_t,
\qquad
D_{jst}=s_{jst}-\operatorname{median}_{k\ne j}\{s_{kst}\},
\label{eq:signal}
\end{equation}
where \(i_{jst}\) is the lending rate, \(r_t\) the policy rate, \(s_{jst}\)
the spread over policy, and \(D_{jst}\) the spread relative to peer banks in
the same segment and month. Thus \(D_{jst}=0.02\) means two percentage points
above the peer median when rates are decimals. The observed outcome is
\begin{equation}
\operatorname{NPL}_{js,t+h}
=\frac{\text{non-performing loan balance}_{js,t+h}}
       {\text{total outstanding loan balance}_{js,t+h}},
\label{eq:npl}
\end{equation}
using the data provider's loan-classification rule.

\section{Agent problem and economic mechanism}

The model suppresses \(j,s,t\) to make the mechanism readable. Borrower
opacity is \(\kappa\in[0,1]\): a larger value means applicants are harder to
compare using readily available information. The timing is: (1) a segment
with \(\kappa\) is given; (2) the bank posts the relative spread \(D\); (3)
price sorting forms an applicant pool; (4) the bank observes that pool and
chooses screening effort \(e\); and (5) repayment or default is realized.

Before screening, the applicant pool has default probability
\begin{equation}
q^0(D,\kappa)=\bar q+a(\kappa)D,
\qquad \bar q\in(0,1),\quad a(\kappa)>0.
\label{eq:pool}
\end{equation}
Here \(\bar q\) is baseline default risk and \(a(\kappa)\) is the sensitivity
of pool quality to relative price. The positive sign captures the selection
channel: a higher relative price loses more price-sensitive applicants and
worsens the unscreened pool. Screening effort \(e\in[0,\bar e]\) changes risk
to
\begin{equation}
q(D,e,\kappa)=q^0(D,\kappa)-g(\kappa)e,
\qquad g(\kappa)>0,
\label{eq:quality}
\end{equation}
where \(g(\kappa)\) is the reduction in default probability produced by one
unit of effort. Effort costs \(c(\kappa)e^2/2\), with \(c(\kappa)>0\).

Let \(R(D)=\bar R+D>0\) be the payoff from a performing loan,
\(F\) funding and operating cost independent of effort, and
\(L(D,\kappa)>0\) potential loan volume, with \(L_D<0\). The bank solves
\begin{equation}
\max_{D\in\mathcal D} L(D,\kappa)
\left\{\max_{e\in[0,\bar e]}
[1-q^0(D,\kappa)+g(\kappa)e](\bar R+D)
-F-\frac{c(\kappa)}2e^2\right\},
\label{eq:problem}
\end{equation}
where the feasible price set
\(\mathcal D=[\underline D,\overline D]\) is compact. Continuity and compactness
ensure a first-stage price solution. The topic result studies the continuation
quality map \(q^*(D,\kappa)\) for any observed relative price, including an
equilibrium price or variation generated by bank-specific markup and cost
shocks. A particular demand system would be needed for a closed-form \(D^*\),
but not for the conditional forecasting mechanism.

\begin{table}[t]
\centering\small
\begin{tabular}{>{\(}l<{\)} p{0.74\textwidth}}
\toprule
\text{Symbol} & \textbf{Meaning} \\
\midrule
D & relative spread chosen in stage one \\
\kappa & borrower-segment opacity \\
q^0,\ q & default probability before and after screening \\
e,\ \bar e & screening effort and its upper bound \\
a & price-sorting sensitivity of the applicant pool \\
g,\ c & quality gain per effort unit and screening-cost shifter \\
\bar R+D,\ F & performing-loan payoff and effort-independent cost \\
L & potential loan volume \\
\bottomrule
\end{tabular}
\caption{Primitive objects in the model.}
\label{tab:notation}
\end{table}

\section{Expected first-order condition and result}

Fix \(D\). The marginal benefit of effort is the quality gain times the payoff
from performance, \(g(\kappa)(\bar R+D)\); its marginal cost is
\(c(\kappa)e\). The first- and second-order conditions are
\begin{equation}
g(\kappa)(\bar R+D)-c(\kappa)e=0,
\qquad \frac{\partial^2\pi}{\partial e^2}=-c(\kappa)<0.
\label{eq:foc}
\end{equation}
For an interior solution,
\begin{equation}
e^*(D,\kappa)=\frac{g(\kappa)}{c(\kappa)}(\bar R+D).
\label{eq:effort}
\end{equation}
A higher \(D\) raises the payoff from avoiding default and therefore raises
screening effort. Substituting \eqref{eq:effort} into \eqref{eq:quality} gives
\begin{equation}
q^*(D,\kappa)=\bar q+a(\kappa)D
-h(\kappa)(\bar R+D),
\qquad
h(\kappa)\equiv\frac{g(\kappa)^2}{c(\kappa)}.
\label{eq:solution}
\end{equation}
The object \(h\) is \emph{screening leverage}: the improvement in quality
created by the optimal screening response to a marginally larger payoff. It is
not the effort level itself. The central decomposition is
\begin{equation}
\boxed{\beta(\kappa)\equiv\frac{\partial q^*}{\partial D}
=\underbrace{a(\kappa)}_{\text{worse selection}}
-\underbrace{h(\kappa)}_{\text{stronger screening}}.}
\label{eq:beta}
\end{equation}

\textbf{Candidate Proposition 1 (opacity threshold).} Suppose: (i) the
solution is interior and \(q^*\in(0,1)\); (ii) \(a\) and \(h\) are continuous;
(iii) \(a'(\kappa)<0\); (iv) \(h'(\kappa)>0\); and (v)
\(a(0)>h(0)\) while \(a(1)<h(1)\). Then a unique
\(\kappa^*\in(0,1)\) solves \(a(\kappa^*)=h(\kappa^*)\), and
\begin{equation}
\beta(\kappa)
\begin{cases}
>0,&\kappa<\kappa^*,\\
=0,&\kappa=\kappa^*,\\
<0,&\kappa>\kappa^*.
\end{cases}
\qquad
\frac{\partial^2q^*}{\partial D\partial\kappa}=a'-h'<0.
\label{eq:threshold}
\end{equation}
Continuity and endpoint crossing give existence by the intermediate value
theorem; \(a'-h'<0\) gives uniqueness. The key economic restriction is exposed
by logarithmic differentiation:
\begin{equation}
\frac{h'}h=2\frac{g'}g-\frac{c'}c.
\label{eq:keycondition}
\end{equation}
Therefore, higher screening cost alone cannot create the negative slope: it
reduces \(h\). Opacity must raise the marginal information gain \(g\) fast
enough to dominate any rise in \(c\). This is a falsifiable sufficient
condition, not an implication built into the definition of opacity.

The theory's \(q^*_{jst}\) is default probability for newly originated credit,
whereas \(\operatorname{NPL}_{js,t+h}\) is a later stock ratio. I connect them
through the explicit forecasting map
\begin{equation}
\operatorname{NPL}_{js,t+h}=\alpha_h+\rho_hq^*_{jst}+u_{js,t+h},
\qquad \rho_h>0.
\label{eq:mapping}
\end{equation}
If this monotone map is supported, the predictive slope has the sign of
\(\rho_h\beta(\kappa)\), hence the sign of \(\beta\). Equation
\eqref{eq:mapping} is a bridge to measurement, not a causal structural
equation.

\section{Why this is not already solved}

The ingredients exist separately. \citet{varian1980} explains price dispersion
using informed and uninformed buyers, but has neither screening nor default.
\citet{broecker1990} combines rate competition with imperfect credit tests, but
does not make testing a continuous effort choice. \citet{ruckes2004}
endogenizes screening intensity as economic prospects change, but does not
characterize the sign of a relative-price signal across opacity segments.
\citet{dellariccia2006} shows how information asymmetry across banks selects
screening or pooling regimes, rather than deriving a continuous
\(\partial q^*/\partial D\) decomposition. The narrow contribution here is to
combine relative-price sorting with endogenous screening and identify the
condition for a unique opacity threshold.

\section{Plan, validation, and risks}

I will (1) write a full proof of Proposition 1 and a robustness result for a
general increasing payoff \(R(D)\); (2) simulate the complete admissible
parameter region instead of reporting only one calibration; (3) test whether
\(D\) improves out-of-sample NPL-ratio forecasts by opacity group and horizon;
and (4) formalize the threshold theorem in Lean after its assumptions are
frozen. Current code symbolically checks the FOC, SOC, and derivative and
numerically checks crossing, monotonicity, interior effort, and probability
bounds.

The main risk is that \(h'(\kappa)>0\) fails. I will then report the three
regions \(a>h\), \(a=h\), and \(a<h\) without forcing a threshold. A second
risk is endogenous pricing. The empirical fallback is explicitly predictive:
bank and segment-by-month controls, rolling out-of-sample evaluation, and no
causal interpretation of the price coefficient.

\small
\bibliographystyle{plainnat}
\bibliography{../paper/references}

\end{document}
